
\renewcommand\thesection{\arabic{section}}
\renewcommand{\contentsname}{Cuprins}
\renewcommand{\figurename}{Figura}
\renewcommand{\tablename}{Tabel}

\setcounter{tocdepth}{6}
\setcounter{secnumdepth}{6}


% \usepackage[backend=bibtex,style=numeric,citestyle=numeric,backref=true,sorting=none]{biblatex}




\titleformat{\section}{\normalfont\Large\justifyheading}{\thesection}{18pt}{}



\definecolor{codegreen}{rgb}{0,0.6,0}
\definecolor{codegray}{rgb}{0.5,0.5,0.5}
\definecolor{codepurple}{rgb}{0.58,0,0.82}
\definecolor{backcolour}{rgb}{0.96,0.96,0.96}

\lstdefinestyle{mystyle}{
	backgroundcolor=\color{backcolour},   
	commentstyle=\color{codegreen},
	keywordstyle=\color{blue},
	numberstyle=\tiny\color{codegray},
	stringstyle=\color{codepurple},
	basicstyle=\fontsize{10}{10}\selectfont\ttfamily,
	breakatwhitespace=false,         
	breaklines=true,                 
	captionpos=b,                    
	keepspaces=true,                 
	numbers=left,                    
	numbersep=2pt,                  
	showspaces=false,                
	showstringspaces=false,
	showtabs=false,                  
	tabsize=2
}

\lstset{style=mystyle}




\setmainfont{UTSans-Regular.otf}


	


\newpage



\newpage
\pagenumbering{arabic}
\renewcommand{\bibname}{Bibliografie}
\renewcommand\lstlistlistingname{\normalsize{CODURI SURSĂ}}
\renewcommand{\lstlistingname}{\normalsize{Secvență de Cod}}
\renewcommand{\listfigurename}{FIGURI}
\renewcommand{\cftloftitlefont}{\normalsize}
\renewcommand{\cftfigfont}{\normalsize}
\renewcommand{\cftfigpagefont}{\normalsize}
\renewcommand{\listtablename}{TABELE}
\renewcommand{\cftlottitlefont}{\normalsize}
\renewcommand{\cfttabfont}{\normalsize}
\renewcommand{\cfttabpagefont}{\normalsize}
 \clearpage

% % trebuie completata manual

\clearpage




\setcounter{page}{4}
\section{Introducere}
În ultimul timp colaborarea dintre oameni și roboți în mediile industriale a prins tracțiune, în special în procesele de asamblare și dezasamblare de mare precizie. Creșterea complexității mediului de lucru industrial, necesită metodologii de antrenare avansate, prin integrarea expertizei umane cu cea a asistenței din partea roboților. Dar această instruire necesită resurse ridicate și implică riscuri de accidentare.

Tehnologia VR reprezintă o soluție modernă care permite simularea și controlul acestor roboți, oferindu-le utilizatorilor posibilitatea de a învăța cum să-i asambleze, dezasambleze, controleze, limitele lor și scenariile de utilizare fără a fi necesar interacțiunea fizică cu cei reali. Un sistem de antrenare VR poate reduce costurile, crește eficiența în antrenare și de a îmbunătăți nivelul de pregătire fără a afecta productivitatea reală.

În prezent au fost propuse inițiative promițătoare în antrenarea folosind  realitate virtuală de antrenare a utilizatori umani de a folosi roboți cum sunt: VR Co-lab \cite{ref1}, sistem de antrenare în realitatea virtuală în mediul industrial \cite{ref2} și a unui sistem de antrenare în realitatea virtuală pentru roboți industriali \cite{ref3}.

Această lucrare își propune realizarea unui sistem interactiv de antrenare folosind realitatea virtuală pentru mediul industrial care lucrează cu roboți care să faciliteze antrenarea utilizatorilor umani neexperimentați la un mediul complet de antrenare în folosirea roboților fără implicarea de riscuri.

\subsection{Obiectivele și scopul lucrării}
Scopul lucrării este de a crea un mediu de simulare interactivă , simplu și sigur pentru antrenarea în folosirea roboților care este destinată utilizatorilor neexperimentați.

De asemenea acest sistem urmărește ca procesul de antrenare a utilizatorilor umani, care nu au interacționat niciodată sau rar cu un astfel de sistem să fie cât mai simplă, intuitivă astfel încât timpul petrecut în antrenare să fie cât mai mic și nivelul de pregătire dobândit să fie cât mai mare.

Nu în ultimul rând prin folosirea acestui mediu de antrenare procesul de antrenare nu implică riscurile aferente, dacă antrenarea implica roboții reali.



Această lucrare prezintă următoarea structură:capitolul 2 prezintă tehnologiile și fundamentele teoretice care stau la baza realizării acestui proiect, capitolul 3 prezintă echipamentele folosite, implementarea, arhitectura aplicației în  realitatea virtuală pentru și modul în care robotul a fost transpus în acest mediu, iar capitolul 4 cuprinde testele realizate și rezultatele în urma testelor, modul în care au fost îndeplinite obiectivele, utilitatea acestui proiect și direcțiile viitoare.
     



% \nocite{*}


